\SourceLeaf{19}
in rechteckige Elemente zerlegt. Ein solches Element liegt zwischen den Punkten mit den Koordinaten
\[
\begin{array}{ll}
p,&q\\
p+\mathrm{d}p,&q\\
p,&q+\mathrm{d}q\\
p+\mathrm{d}p,&q+\mathrm{d}q
\end{array}
\]
Sein Flächeninhalt ist $=\mathrm{d}p.\mathrm{d}q$, und ihm entspricht auf der Körperoberfläche ein parallelogrammförmiges Element zwischen vier Punkten, deren Koordinaten der Reihe nach lauten
\[
\begin{array}{rlll}
\mathrm{I.}&x,&y,&z\\
\mathrm{II.}&x+\lambda\mathrm{d}p,&y+\mu\mathrm{d}p,&z+\nu\mathrm{d}p\\
\mathrm{III.}&x+\lambda'\mathrm{d}q,&y+\mu'\mathrm{d}q,&z+\nu'\mathrm{d}q\\
\mathrm{IV.}&x+\lambda\mathrm{d}p+\lambda'\mathrm{d}q,&y+\mu\mathrm{d}p+\mu'\mathrm{d}q,&z+\nu\mathrm{d}p+\nu'\mathrm{d}q
\end{array}
\]
wenn wir annehmen, dass Folgendes gilt
\[
\begin{aligned}
\mathrm{d}x&=\lambda\mathrm{d}p+\lambda'\mathrm{d}q\\
\mathrm{d}y&=\mu\mathrm{d}p+\mu'\mathrm{d}q\\
\mathrm{d}z&=\nu\mathrm{d}p+\nu'\mathrm{d}q
\end{aligned}
\]
Die Projektionen dieser Fläche, deren Inhalt wir $=\mathrm{d}s$ setzen, auf die drei zu den Koordinatenachsen $x,y,z$ senkrechten Ebenen ergeben sich leicht der Reihe nach zu $=$
\[
\begin{aligned}
&\pm(\mu\nu'-\nu\mu')\mathrm{d}p.\mathrm{d}q\\
&\pm(\nu\lambda'-\lambda\nu')\mathrm{d}p.\mathrm{d}q\\
&\pm(\lambda\mu'-\mu\lambda')\mathrm{d}p.\mathrm{d}q
\end{aligned}
\]
Nach einem wohlbekannten Lehrsatz ergibt sich daraus für den Flächeninhalt des Elements selbst
\[
=\mathrm{d}p.\mathrm{d}q.\BareRoot{((\mu\nu'-\nu\mu')^2+(\nu\lambda'-\lambda\nu')^2+(\lambda\mu'-\mu\lambda')^2)}
\]
\par\endgroup
\par
Daraus folgt, dass sich jedes der in unseren sechs Lehrsätzen angegebenen Integrale auf die Form $\int S\mathrm{d}p.\mathrm{d}q$ bringen lässt, wobei $S$ explizit oder implizit eine Funktion der beiden Unbestimmten $p,q$ ist. Die Integration ist entweder über die gesamte unendliche Ebene oder über den Teil der Ebene zu erstrecken, der gewissermaßen die ganze Oberfläche unseres Körpers darstellt. Die Integration selbst lässt sich bald mit diesem, bald mit jenem Kunstgriff durchführen. Allgemeine Regeln dafür lassen sich nicht angeben.
